\subsection{Durability of the innercode}
\subsubsection{Primitives}

Suppose a system with $N$ total nodes. There exists $F$ byzantine nodes, where $F = \frac{N}{3}$ and the number of honest nodes is $N-F$. 

Out of these $N$ nodes, a group of $n$ nodes is selected (without replacement, since only unique identities are considered), at any $T\in \mathbb{Z}^{+}$ there needs to be very high probability that:
\begin{enumerate}
    \item The initial state of $n$ at time $T=0$, there exists at least a certain threshold of honest nodes $k$, with overwhelming probability. Suppose then max number of byzantine nodes are $b = n-k$ 
    \item Given a churn rate that is expressed with the Poisson distribution, there must still always exist $k$ honest nodes after some time $T=t,t\in \mathbb{Z}^{+}$. 
\end{enumerate}

$N$ and $n$ are parameters to this algorithm, where $k$ would exist as a fixed constant. 

\subsubsection{Intitial chunk state durability}

We say a state is valid if $Pr(b>n-k) < \varepsilon = 2^{-128}$. Since we know that group selection is expressed with the hypergeometric distribution, there are two way to show this:
\begin{enumerate}
    \item Utilizing the generalized hypergeometric function or summation of PMF to derive the CDF. 
    \item Utilizing chebyshev's inequality to draw a bound. 
\end{enumerate}


\begin{gather}
    P(b > n-k) = 1- \sum_{i=0}^{n-k}\frac{\binom{\frac{N}{3}}{i}\binom{\frac{2N}{3}}{n-i}}{\binom{N}{n}} \\
    \leq 
    \frac{n\frac{F}{N}\frac{N-F}{N}\frac{N-n}{N-1}}{(n(1+\frac{F}{N})-k)^2} \\
    \leq 
    \frac{n(N-n)}{(4n-3k)^2(N-1)} |_{F=\frac{N}{3}}\\
    \leq \varepsilon
\end{gather}

However, Although the intitial state might be valid, we cannot say future states will be. The validity of any state at $T=t$ is evaluated in the next section. 

One thing to note is that, if our $N$ shrinks, $Pr(b>n-k)$ actually decreases. This can be seen in $(1)$, as $N$ decreases, each summation step increases slightly, making the inverse of the probability sum smaller. 

\subsubsection{Bounded state chunk durability}

\begin{lemma}
    \label{lemma 1}
    For a single data object, the probability that all of the groups injectively responsible for the $K+R$ chunks being durable at $T=t$ can be bounded by: \[1-(1-\sum_{T=1}^{t} (\pi P^{T})_{n-k+1})^{K+R} \leq \varepsilon\] Where $\pi$ is the initial state matrix and $\Theta$ is a stochastic matrix. 
\end{lemma}

We analyze our system as a Continuous-time Markov chain (CMTC). Expressing our initial group state as a $(n-k+1 \times 1)$ probability vector:

\begin{gather}
    I = \begin{pmatrix}
            Pr(B=0) \\ 
            Pr(B=1)\\
            \vdots \\
            Pr(B \geq n-k+1)
            \end{pmatrix}, 
    Pr(B = b) = \frac{\binom{F}{b}\binom{N-F}{n-b}}{\binom{N}{n}}
\end{gather}

We express the probability of reaching any state $T=t$ as $t$ exponentiations of a $(n-k+1 \times n-k+1)$ stochastic matrix $\Theta = (\theta_{i,j})$ on this intitial probability vector.
$\theta_{i,j}$ denotes the probability a group with $i$ byzantine nodes at $T=t$ will have $j$ byzantine nodes at $T=t+1$. 

Since we take our group as invalid in perpetuity once its reached an invalid state; These states are referred to as \textit{absorbing states}. Any other state is considered as an \textit{transient states}. Our matrix is only $(n-k+1 \times n-k+1)$ instead of $(n \times n)$. Furthemore, $\theta_{n-k+1,n-k+1} = 1 \wedge \forall g\in[0,n-k], \theta_{n-k+1,g} = 0$. This handles all transitions from \textit{absorbing} to \textit{absorbing states}. 

To construct the remainder of our matrix, we will have to consider the possible state transitions of each originating state. Suppose for some $\lambda(N-F)$ , let $C$ be the random variable denoting the number of honest nodes lost due to the churn rate:
\begin{gather}
    Pr(C=c) = \frac{\lambda(N-F)^c e^{-c}}{c!}
\end{gather}

Furthermore, our system has a fixed eviction rate $\Upsilon$, which denotes the number nodes from all nodes that are selected randomly to be removed from the group. 

\begin{gather}
    \forall i,j \in [0, n-k],\\
    \theta_{i,j} = 
    \Biggl( 
    \sum_{c=0}^{n-i-k} 
    \frac{\lambda(N-F)^c e^{-c}}{c!} \\
    %This sum is the probabilities of all the churn results that will result in a durable state. Ie: c <- number of honest nodes that leave
    \Biggl( 
    \sum_{\substack{\upsilon=0\\(j-i+\Upsilon-\upsilon) \geq 0\\(n-k-c)\geq \Upsilon} }^{min(n-i-k-c,\Upsilon)} 
    \frac{\binom{i}{\Upsilon-\upsilon}\binom{n-i-c}{\upsilon}}{\binom{n-c}{\Upsilon}} \\
    %At this point, (n-i-c) honest nodes are left. Therefore, i can evict a range of [0 honest nodes, min(n-i-c,Upsilon)] honest nodes. 
    %in other words, if i evict v honest nodes, im actually evicting Upsilon-v honest nodes. 
    %NOTE: if j-(i-(Upsilon-v))  is negative, we dont count, we just simply ignore and go to a churn rate that matters. 
    \frac{\binom{F-i-(\Upsilon-\upsilon)}{j-i+\Upsilon-\upsilon}\binom{N-(n-c-\Upsilon)-(F-i-(\Upsilon-\upsilon))}{n-c-\Upsilon-(j-i+\Upsilon-\upsilon)}}{\binom{N-(n-c-\Upsilon)}{n-c-\Upsilon}}
    %now there is a remainder of i-(Upsilon-v) byzantine nodes. 
    %that means we need to find j-(i-(Upsilon-v)) or (j-i+\Upsilon-\upsilon) byzantine nodes
    \Biggl)\Biggl)
\end{gather}

Since this only considers $\forall i,j \in [0, n-k]$, it only concerns \textit{transient} to \textit{transient state} transitions. The first summation (8) considers all the possible valid number honest nodes churned, such that at least $k$ honest nodes remain. 
The second summation (9) considers the eviction rate; It first considers the amount of honest nodes that can be evicted $min(n-i-k-c,\Upsilon)$. For example, if the amount of honest nodes after churning is $k$, then obviously only byzantine nodes can be evicted if the next state is a transient state. This information is used to calculate how many byzantine nodes will be evicted ($\Upsilon - \upsilon$).

The condition in the second summation: $(n-k-c)\geq \Upsilon$ checks if there exists sufficient nodes to be evicted, and $(j-i+\Upsilon-\upsilon) \geq 0$ considers if the next state can be reached after churning and eviction. 

The first fraction in the second summation (9) thus considers how many byzantine nodes can be evicted. The second fraction (10) then considers the probability of adding back exactly $(j-i+\Upsilon-\upsilon)$ byzantine nodes back into the group to acheive the desired transient state. 

Lastly, we consider transitions from \textit{transient} to \textit{absorbing states}:

\begin{gather}
    % we consider 3 main cases:
    % 1. We churned too much
    % 2. we churned okay, but evicted too much
     \forall i\in[0, n-k], j = n-k+1, \\
     \theta_{i,j} = 
     1- \sum_{g = 0}^{n-k} \theta_{i,g}
\end{gather}

The probability that a group has entered an absorbing state is then expressed as:

\begin{gather}
    \sum_{T=1}^{t} (\pi \Theta^{T})_{n-k+1}
\end{gather}

Given enough time, the probability of reaching any state is $1$, therefore, that is to say there is no durability at $T=\infty$. 

However, we can guarantee at some $T=t$, the probability the group has reached an invalid state is still negligable. By extension, it also shows that the probability the group has reached an invalid state for $T=t'|t'\in \mathbb{Z}^{+} \wedge t'<t$ is negligable too. 

Suppose a data object is made up of $K+R$ chunks, the probability of some group/s reaching an  \textit{absorbing} state is:

\begin{gather}
    1-(1-\sum_{T=1}^{t} (\pi \Theta^{T})_{n-k+1})^{K+R} \leq \varepsilon \qquad\blacksquare
\end{gather}


\subsection{Durability in the event of targetted attacks}

\subsubsection{Primitives}
Suppose in our system of $\Omega$ data objects, each object is made up of $K+R$ chunks; where only a minimum of $K$ chunks is necessary to reconstuct the original data object. $R$ thus represents the redundancy.

We present an adversary that can arbitrarily remove an upperbound of $\phi$ nodes from the system; Effectively churning them pre-maturely.
Furthermore, the adversary has a complete transparent view on the group composition for every group.
This begs the question: how many groups can they arbitrarily push into the absorbing state, and what is the probability a data object will be lost with non-negligable odds. We consider a successful attack on a group if the attacker can lead it to an absorbing state, and we consider a sucessful attack on the system if any data objects are lost as the result of the attacker's actions. 

There are two cases to consider, the latter being an extension of the former. 

\subsubsection{Suppose each phyiscal node only holds a single fragment}
\begin{lemma} \label{lemma2}
    The probability of an attacker successfully attacking a data object is bounded by:
    \[
        1-
    \left(
    1-\prod_{i=1}^{R}\frac{K+R-i}{\Omega(K+R)-i}
    \right)^{\binom{\Phi}{R+1}} \leq \varepsilon
        \]
\end{lemma}

We first present the naive scenario in which each node holds on to a single fragment. That means attacking a single node will at most lead to a single fragment from being lost. Another way to visualize this is that one node can only exist in a single group. 

On average, each group should have an average of $\frac{n}{3}$ honest nodes (per the hypergeometric distribution). This means the attacker will have to attack $\frac{n}{3} - k+1$ nodes on average to lead the group into an absorbing state, or $\Phi= \lfloor\frac{\phi}{\frac{n}{3}-k+1} \rfloor$ groups on average. 


However, since the attacker has a perfect view on the state of the system, it does not attack the system randomly. Instead, it can pick the most opportune time to maximize its attack vector. If suppose $\geq \phi$ groups have exactly $k$ honest nodes in their groups the attacker will at worst successfully attack $\Phi  = \phi$ groups. 

Upperbound on the Probability of successfully attacking a data object is thus (suppose $\Phi \geq R+1$) :
\begin{gather}
    1-
    \left(
    1-\prod_{i=1}^{R}\frac{K+R-i}{\Omega(K+R)-i}
    \right)^{\binom{\Phi}{R+1}}  \leq \varepsilon
    \qquad\blacksquare
\end{gather}

This is an extension of the birthday attack problem. 

\subsubsection{Suppose each phyiscal node only holds multiple fragment}
Since we are drawing an upperbound on the attacker's success probability, we simply have to consider the very worst case. 
If each node can possess $g,g\in \mathbb{Z}^{+}$ fragmemts, attacking a single node might lead to $g$ groups losing a single honest node each. 
Therefore, we simply change (14) to:
\begin{gather}
    1-
    \left(
    1-\prod_{i=1}^{R}\frac{K+R-i}{\Omega(K+R)-i}
    \right)^{\binom{\Phi*g}{R+1}} \leq \varepsilon
\end{gather}
This should be the upperbound on the Probability of successfully attacking a data object. 